% ============================================================
%  参考文献
%  课程要求：限于优质的中英文期刊。
%  政府文件、国家标准、法律法规不走此表，作为正文出处标注。
%  英文条目经 Crossref 按 DOI 取权威元数据；正文采用作者—年份制。
%  本表按作者字母／拼音顺序排列。
% ============================================================
\section*{参考文献}
\addcontentsline{toc}{section}{参考文献}

\begin{list}{}{\setlength{\leftmargin}{2.2em}\setlength{\itemindent}{-2.2em}\setlength{\itemsep}{0.35em}\raggedright\sloppy}

\subsection*{中文文献}

\item 安体富, 樊勇. 对我国消费税改革的评析与展望. \emph{税务研究}, 2006(12): 35--37.
\item 陈欣, 等. 我国重点城市春节期间空气质量变化特征. \emph{统计与信息论坛}, 2014, 29(12): 57--62.
\item 褚睿刚. 消费税“绿化”功能的检思与超越——兼论污染产品税的创立. \emph{经济法学评论}, 2016, 16(2): 279--301.
\item 李曼, 吴瑶, 田耀华, 等. PM\textsubscript{2.5} 短期暴露对我国死亡和急诊量影响的 Meta 分析. \emph{中华流行病学杂志}, 2018, 39(10): 1394--1401.
\item 陆卜梵, 刘蓓蓓. 京津冀地区烟花爆竹禁放政策对空气质量的影响. \emph{环境科学学报}, 2022, 42(8): 26--35.
\item 那艺诺, 李艺, 王丹, 等. 基于支付意愿法中国统计生命价值系统评价. \emph{中国公共卫生}, 2025, 41(2): 239--245.
\item 庞凤喜, 王绿荫. 消费税改革的目标定位及制度优化分析. \emph{税务研究}, 2020(1): 44--49.
\item 任国保, 何俊杰. 推进鞭炮焰火行业高质量发展的税收思考——以湖南鞭炮焰火行业为例. \emph{湖南税务高等专科学校学报}, 2020, 33(6): 32--36.
\item 田璐, 金曼, 佟俊旺. 应用 Meta 分析研究我国 PM\textsubscript{2.5} 污染与居民每日死亡的关系. \emph{环境与职业医学}, 2015, 32(11): 1013--1018.

\subsection*{英文文献}

\item Beiser-McGrath, L. F., Bernauer, T., \& Prakash, A. (2023). Command and control or market-based instruments? Public support for policies to address vehicular pollution in Beijing and New Delhi. \emph{Environmental Politics}, 32(4), 586--618. DOI: 10.1080/09644016.2022.2113608.
\item Buchanan, J. M., \& Tullock, G. (1975). Polluters' profits and political response: Direct controls versus taxes. \emph{American Economic Review}, 65(1), 139--147.
\item Callaway, B., \& Sant'Anna, P. H. C. (2021). Difference-in-differences with multiple time periods. \emph{Journal of Econometrics}, 225(2), 200--230. DOI: 10.1016/j.jeconom.2020.12.001.
\item Camilleri, R., \& Vella, A. J. (2010). Effect of fireworks on ambient air quality in Malta. \emph{Atmospheric Environment}, 44(35), 4521--4527. DOI: 10.1016/j.atmosenv.2010.07.057.
\item Cao, X., et al. (2018). Review on physicochemical properties of pollutants released from fireworks. \emph{Environmental Reviews}, 26(2), 133--155. DOI: 10.1139/er-2017-0063.
\item Chen, S., Jiang, L., Liu, W., \& Song, H. (2022). Fireworks regulation, air pollution, and public health: Evidence from China. \emph{Regional Science and Urban Economics}, 92, 103722. DOI: 10.1016/j.regsciurbeco.2021.103722.
\item Croteau, G. A., Dills, R. L., Beaudreau, M., \& Davis, M. (2010). Emission factors and exposures from ground-level pyrotechnics. \emph{Atmospheric Environment}, 44(27), 3295--3303. DOI: 10.1016/j.atmosenv.2010.05.048.
\item de Chaisemartin, C., \& D'Haultf\oe uille, X. (2020). Two-way fixed effects estimators with heterogeneous treatment effects. \emph{American Economic Review}, 110(9), 2964--2996. DOI: 10.1257/aer.20181169.
\item Harrington, W., Morgenstern, R. D., \& Sterner, T. (Eds.) (2004). \emph{Choosing Environmental Policy: Comparing Instruments and Outcomes in the United States and Europe}. Washington, DC: RFF Press.
\item Jiang, Q., Sun, Y. L., Wang, Z., \& Yin, Y. (2015). Aerosol composition and sources during the Chinese Spring Festival: Fireworks, secondary aerosol, and holiday effects. \emph{Atmospheric Chemistry and Physics}, 15(11), 6023--6034. DOI: 10.5194/acp-15-6023-2015.
\item Keller, K., \& Schragen, C. (2021). Determination of particulate matter emission factors of common pyrotechnic articles. \emph{Propellants, Explosives, Pyrotechnics}, 46(5), 825--842. DOI: 10.1002/prep.202000292.
\item Li, C., Wang, J., Yan, B., et al. (2026). Dynamic correction of E/C/$\gamma$ parameters and construction of human health quantification system for air pollution damage assessment based on virtual governance cost method. \emph{Journal of Cleaner Production}, 558, 148283. DOI: 10.1016/j.jclepro.2026.148283.
\item Li, Y., Li, R., Wang, S., \& Wei, W. (2025). Fireworks celebrations and air pollution: Evidence from Chinese Lunar New Year. \emph{Journal of Environmental Planning and Management}, 68(3), 640--660. DOI: 10.1080/09640568.2023.2268268.
\item Liu, C., Chen, R., Sera, F., et al. (2019). Ambient particulate air pollution and daily mortality in 652 cities. \emph{New England Journal of Medicine}, 381(8), 705--715. DOI: 10.1056/NEJMoa1817364.
\item Oates, W. E., \& Baumol, W. J. (1975). The instruments for environmental policy. In E. S. Mills (Ed.), \emph{Economic Analysis of Environmental Problems}. Cambridge, MA: NBER.
\item Pang, N., Gao, J., Zhao, P., Wang, Y., Xu, Z., \& Chai, F. (2021). The impact of fireworks control on air quality in four Northern Chinese cities during the Spring Festival. \emph{Atmospheric Environment}, 244, 117958. DOI: 10.1016/j.atmosenv.2020.117958.
\item Pang, N., Jiang, B., \& Zheng, Y. (2023). The impact of fireworks burning on air quality and their health effects in China during Spring Festivals of 2015--2022. \emph{Atmospheric Pollution Research}, 14(11), 101888. DOI: 10.1016/j.apr.2023.101888.
\item Tang, A., Tang, A. Y. F., \& Peng, P. Z. J. (2012). Recycling of global consumer fireworks. \emph{Journal of Pyrotechnics}, 31, 22--32.
\item Wang, Y., et al. (2024). Valuing mortality risk reductions in a fast-developing society: A meta-analysis of stated preference studies in China from 1998 to 2019. \emph{Social Science \& Medicine}, 363, 117471. DOI: 10.1016/j.socscimed.2024.117471.
\item Weitzman, M. L. (1974). Prices vs.\ quantities. \emph{The Review of Economic Studies}, 41(4), 477--491. DOI: 10.2307/2296698.
\item Yang, J., Wei, C., Shen, R., \& Ji, T. (2026). Targeted environmental policy and public welfare: Evidence from fireworks deregulation in China. \emph{China Economic Review}, 96, 102668. DOI: 10.1016/j.chieco.2026.102668.
\item Zhang, X., Shen, H., Li, T., \& Zhang, L. (2020). The effects of fireworks discharge on atmospheric PM\textsubscript{2.5} concentration in the Chinese Lunar New Year. \emph{International Journal of Environmental Research and Public Health}, 17(24), 9333. DOI: 10.3390/ijerph17249333.

\end{list}
